\chapter[Q5 Is the Random Process Ergodic]{Q5 Is the Random Process Ergodic}
\label{cp:q5-is-the-random-process-ergodic}

To determine if a random process is ergodic, we need to satisfy the following
conditions as learned in \cite{DigitalCommunicationsLecture4}, statistics across time are the same
as their statistics across the ensemble, which means:
\begin{enumerate}
    \item \begin{equation}
              \label{eq:moment1}
              m(t) = E(X(t)) = \langle X(t) \rangle
          \end{equation}

    \item \begin{equation}
              \label{eq:moment2}
              R_x(\tau) = E[X(t) X(t + \tau)] = \langle X(t) X(t + \tau) \rangle
          \end{equation}
\end{enumerate}

Where $E$ is the expectation across the ensemble and $\langle \cdot \rangle$ is
the time average.

We can use the following equations to get the time average of the mean and
autocorrelation
\begin{equation}
    \langle x(t) \rangle = \lim_{T\to\infty} \frac{1}{2T} \int_{-T}^{T} x(t) \, dt
\end{equation}

\begin{equation}
    \langle x(t) x(t + \tau) \rangle = \lim_{T\to\infty} \frac{1}{2T} \int_{-T}^{T} x(t)x(t + \tau) \, dt
\end{equation}

We have already calculated the mean and autocorrelation across the ensemble in
the previous questions, we also graphed across time and acroess realizations,
the two conditions in \autoref{eq:moment1} and \autoref{eq:moment2} are
satisfied for all three formats as observed from Figures
\autoref{fig:UNIPOLAR_NRZ_MEAN},
\autoref{fig:POLAR_NRZ_MEAN}, 
\autoref{fig:POLAR_RZ_MEAN}, 
\autoref{fig:UNIPOLAR_NRZ_AUTOCORRELATION},
\autoref{fig:POLAR_NRZ_AUTOCORRELATION}, and
\autoref{fig:POLAR_RZ_AUTOCORRELATION}, which means that the
random process is ergodic.
